% 来源：OEBPS/Text/01_31.xhtml（OPF spine 顺序；仅第一册）
\cleardoublepage
\BookEntry[约翰·肯尼思·加尔布雷思《富裕社会》]{约翰·肯尼思·加尔布雷思《富裕社会》}

\BookCover{images/书单17富裕社会.jpg}{[加] 约翰·肯尼思·加尔布雷思著| 赵勇译| 江苏人民出版社2009年版}

\begin{BookSummaryBox}
\textbf{纽约图书馆的世纪丛书之一}
\par
\end{BookSummaryBox}

\subsubsection*{这本书谈什么？}

《富裕社会》由美国自由派经济学家加尔布雷思于50年前写就第一版，后来几经增删和再版。这本书的前几版通篇都在明确表述对通货膨胀的担心。通货膨胀一直对社会基本福利构成威胁。
\par

作者加尔布雷思笔下的“富裕社会”，可以理解为一个经济不断发展的社会，这个社会已经具有消灭贫困和流离失所的力量，只不过为了避免打击那些创造了大量财富的阶层的发展积极性，这样一个“富裕社会”漠视调节社会收入分配的重要性。
\par

在书中，加尔布雷思以“富足社会”与“传统智慧”扩展了经济学的语言，指出了现代富足社会背后隐藏的严重问题，如公共生活与私人生活标准间的资金与行动分配不均，以及环境问题的重要性。
\par

\subsubsection*{芒格为什么谈起这本书？}

2000年11月10日，芒格在慈善圆桌会议早餐会上的讲话，主要内容为美国股市上涨带来的所谓“财富效应”。美国的股票价格对经济产生了多大的影响呢？芒格认为，传统经济学家考虑这个问题所用的数据与现实有很大脱节。除了没有考虑退休金的因素外，也没有把“黑金”（bezzle）的因素考虑在内。
\par

\begin{quotation}
“黑金”这个词跟贪污有关，哈佛大学经济学系教授约翰·肯尼思·加尔布雷思用它指在尚未败露的贪污中得到的金钱。加尔布雷思发现，黑金对消费有非常强烈的刺激作用。毕竟，贪污者花钱更大手大脚，因为他的钱来得更容易，而且他的雇主的支出将会一如既往，因为雇主尚不知道其财产已经被“偷”走了。
\par
\end{quotation}

在芒格看来，加尔布雷思并没有铺开他的洞见，他满足于（在经济学思考上）提供一些牛虻式的刺激。所以芒格进一步发挥了加尔布雷思的“黑金”概念，提出“灰金”“捞灰金”和“灰金客”的概念：
\par

\begin{quotation}
正如凯恩斯指出的，在依靠劳动换取收入的原始经济中，当女裁缝把一件衣服以20美元的价格卖给鞋匠时，鞋匠就少了20美元可以消费，而女裁缝则多了20美元可以用。总消费支出并没有受到lollapalooza效应\footnote{芒格认为，几乎每个系统都受到多种因素的影响，所以要理解这样的系统，就必须熟练运用来自不同学科的多元思维模型。当几个模型联合起来，共同作用一个方向，就得到了芒格所谓的“lollapalooza效应”。——编者注}的影响。但如果政府印刷了另外一张20美元的钞票，用它来买一双鞋，鞋匠多得到了20美元，可是没有人觉得自己的钱变少了。当鞋匠下次再买一件衣服的时候，这个过程就重演了，不会无休止地持续放大，但会产生所谓的凯恩斯乘数效应，这是一种促进消费的lollapalooza效应。同样地，和同等规模的诚实交易相比，尚未败露的贪污得来的钱对消费的刺激效应更大。加尔布雷思是苏格兰人，喜欢深刻地揭示生活的世态炎凉。毕竟，这个苏格兰人还热衷于接受命中注定、无法改变的婴儿诅咒这样的荒唐想法。我们大多数人并不喜欢加尔布雷思的观点。但我们不得不承认，他有关“黑金”的看法基本上是正确的。
\par
\end{quotation}

加尔布雷思出生在加拿大安大略省，从安大略农学院毕业后，到伯克利加利福尼亚大学攻读博士学位。他于1949年到哈佛大学经济学系担任教职，是肯尼迪总统的好友，1961年到1963年担任美国驻印度大使。作为一名经济学家，加尔布雷思持有改良派的价值观，他写了许多描述经济理论如何与现实生活相悖的著作。《富裕社会》便是他的代表作之一。
\par

\subsubsection*{《富裕社会》作者约翰·肯尼思·加尔布雷思如是说}

“我们都认为悲观是智力超群的标志。”
\par

“美当然没有绝对的标准。所以追求美才会如此有趣。”
\par

“如果你什么都不想做，那么你就去开会吧。”
\par

“经济预测的唯一功能是让占星学变得令人尊敬。”
\par

“现代的保守主义者从事的是道德哲学中最古老的活动，也就是说，他们要为自私寻找一个超级漂亮的道德借口。”
\par

“许多美国人死于食物过多，而非死于食物太少。”
\par

“当面临要么改变想法，要么证明无须这么做的选择时，绝大多数人都会忙于寻找证据。那些走运的人总是说他们过得幸福是因为他们的道德很高尚。”
\par

“难道你不明白吗？让一家企业停泊在废话当中，和让它驶入麻烦的思维之海比起来，是一种好得太多、太多的做法。”
\par

“有权势的人总是宁愿拿全副身家去冒险，也不肯舍弃一点利益。”
\par

“在任何大型组织里面，和多数人一起错比一个人对安全得多。”
\par
